\clearpage
\VerdierPhysicalPageAnchor{0176}%
\begingroup\small\setlength{\parskip}{0pt}\setlength{\emergencystretch}{1.5em}
\setlength{\abovedisplayskip}{4pt plus 2pt minus 2pt}%
\setlength{\belowdisplayskip}{4pt plus 2pt minus 2pt}%

\VerdierSemanticAnchor{ch3d-1-2}
\subsection*{1.2. 導来圏の重要な部分圏。}

\VerdierSemanticAnchor{ch3d-1-2-1}
\noindent\textbf{命題 1.2.1.}\quad
\textit{a) 部分圏 $K^{,*}(\mathcal A)\subset K(\mathcal A)$ は右および左
$\operatorname{Ac}(\mathcal A)$-局所化部分圏である（第 II 章，2.3.5）。}

\noindent\textit{b) 部分圏 $K^{+,*}(\mathcal A)\subset K(\mathcal A)$ は
右 $\operatorname{Ac}(\mathcal A)$-局所化部分圏である。}

\noindent\textit{b)\textsuperscript{bis} 部分圏
$K^{-,*}(\mathcal A)\subset K(\mathcal A)$ は左
$\operatorname{Ac}(\mathcal A)$-局所化部分圏である。}

\noindent\textit{c) 部分圏 $K^{b}(\mathcal A)\subset K^{+}(\mathcal A)$ は
左 $\operatorname{Ac}^{+}(\mathcal A)$-局所化部分圏である。}

\noindent\textit{c)\textsuperscript{bis} 部分圏
$K^{b}(\mathcal A)\subset K^{-}(\mathcal A)$ は右
$\operatorname{Ac}^{-}(\mathcal A)$-局所化部分圏である。}

\noindent\textit{d) $K^{,+}(\mathcal A)$ の各対象から
$K^{+}(\mathcal A)$ のある対象への擬同型が存在する。}

\noindent\textit{d)\textsuperscript{bis} $K^{,-}(\mathcal A)$ の各対象は、
$K^{-}(\mathcal A)$ のある対象から擬同型を受ける。}

\noindent\textit{e) $K^{,b}(\mathcal A)$ の各対象から
$K^{+,b}(\mathcal A)$ のある対象への擬同型が存在する。}

\noindent\textit{e)\textsuperscript{bis} $K^{,b}(\mathcal A)$ の各対象は、
$K^{-,b}(\mathcal A)$ のある対象から擬同型を受ける。}

\noindent\textit{f) $K^{+,b}(\mathcal A)$ の各対象は、
$K^{b}(\mathcal A)$ のある対象から擬同型を受ける。}

\noindent\textit{f)\textsuperscript{bis} $K^{-,b}(\mathcal A)$ の各対象から
$K^{b}(\mathcal A)$ のある対象への擬同型が存在する。}

\VerdierSemanticAnchor{ch3d-1-2-1-truncations}
任意の複体 $X$ に対し、$X(n,+\infty)$（それぞれ $X(-\infty,n)$）を複体
\[
X(n,+\infty)=\cdots\longrightarrow0\longrightarrow0
\longrightarrow X^n/\operatorname{Im}(d_X^{n-1})
\longrightarrow X^{n+1}\longrightarrow\cdots
\]
（それぞれ
\[
X(-\infty,n)=\cdots\longrightarrow X^{n-3}\longrightarrow X^{n-2}
\longrightarrow\operatorname{Ker}(d_X^{n-1})
\longrightarrow0\longrightarrow0\longrightarrow\cdots
\]）
とし、$p\colon X\longrightarrow X(n,+\infty)$（それぞれ
$i\colon X(-\infty,n)\longrightarrow X$）を次の複体の射とする：
\[
\begin{aligned}
p^q&=\operatorname{id},&&q>n,\\
p^n&\colon X^n\longrightarrow X^n/\operatorname{Im}(d_X^{n-1})
  &&\text{（標準射）},\\
p^q&=0,&&q\leq n-1.
\end{aligned}
\]
\endgroup
